\section{Analisi del caso: la vulnerabilità di Sign in with Apple del 2020}
\label{sec:analisi-tecnica}

\subsection{Presentazione del caso e contesto della scoperta}
Nell'aprile 2020 il ricercatore di sicurezza indipendente Bhavuk Jain individuò
una grave vulnerabilità nel servizio \emph{Sign in with Apple}, introdotto da
Apple nel 2019 insieme a iOS~13. Il difetto fu segnalato ad Apple attraverso
l'\emph{Apple Security Bounty}~\cite{apple-bounty-2019} e successivamente rilasciato
la notizia al pubblico il 30 maggio 2020, a problema risolto~\cite{jain-apple-2020}.
Alla segnalazione è seguita una ricompensa di 100.000 dollari, e Apple ha dichiarato
di aver analizzato i propri log senza riscontrare abusi o compromissioni di
account legate alla vulnerabilità~\cite{jain-apple-2020,goodin2020}.

Il caso è particolarmente significativo perché il difetto non riguardava
un'applicazione di terze parti, ma direttamente l'\emph{Identity Provider}:
ossia nel componente su cui tutte le applicazioni aderenti al servizio basano
la propria fiducia.

\subsection{Descrizione della vulnerabilità individuata}
Secondo la descrizione del ricercatore, era possibile ottenere da Apple un
token di identità valido per \textbf{qualsiasi indirizzo email} indicato in una
richiesta, anche se tale indirizzo non apparteneva all'utente che stava
effettuando la richiesta~\cite{jain-apple-2020}. In particolare, Apple
rispondeva con un \emph{ID Token} correttamente firmato che attestava
l'appartenenza di quell'email a chi presentava il token. Poiché la firma era
autentica, il token risultava indistinguibile da uno emesso legittimamente.

Il risultato era che un attaccante poteva generare, senza conoscere alcuna
password, un token di identità riferita all'account della vittima.
Il ricercatore precisa che l'attacco funzionava a prescindere dal fatto che
l'indirizzo email indicato fosse associato a un Apple ID valido o meno ~\cite{jain-apple-2020}.

\subsection{Analisi tecnica del meccanismo di attacco}
Il flusso legittimo di \emph{Sign in with Apple} è simile a OAuth~2.0: l'utente
si autentica presso Apple, sceglie se condividere il proprio indirizzo email
reale o un indirizzo anonimo generato da Apple (\emph{relay}) e Apple restituisce
all'applicazione un \emph{ID Token} in formato JWT che contiene, tra gli altri,
l'identificativo dell'utente e l'email~\cite{apple-authentication,jain-apple-2020}.

Nel passaggio in cui veniva richiesto il token, la richiesta inviata ad Apple
includeva un campo con l'indirizzo email. Il problema consisteva nel fatto che
Apple non verificasse che tale valore corrispondesse all'utente effettivamente
autenticato nella sessione, e di conseguenza firmava il token con email della vittima.
L'attacco si articolava quindi in tre passi~\cite{jain-apple-2020}:

\begin{enumerate}
    \item l'attaccante invia ad Apple la richiesta di token indicando l'email
    della vittima;
    \item Apple risponde con un \emph{ID Token} firmato da Apple contenente l'email della
    vittima;
    \item l'attaccante presenta il token all'applicazione di terze parti che
    utilizza \emph{Sign in with Apple}, la quale, se si limita a verificare la
    firma, lo accetta e apre una sessione come se fosse la vittima.
\end{enumerate}

Si tratta di una vulnerabilità di \textbf{controllo dell'autorizzazione lato
IdP}: il sistema autenticava correttamente il richiedente, ma non verificava che
fosse autorizzato a ottenere un'attestazione relativa a quella specifica
identità. Il ricercatore cita come servizi che adottavano il login di Apple
Dropbox, Spotify, Airbnb e Giphy, precisando però di non averne testato
l'effettiva vulnerabilità~\cite{jain-apple-2020}.

\subsection{Credenziali, token e procedure di autenticazione coinvolte}
Gli elementi coinvolti nel processo sono i seguenti~\cite{apple-authentication,apple-verifying-user,openid-connect-core}:

\begin{itemize}
    \item \textbf{Authorization code}: codice temporaneo emesso da Apple al
    termine dell'autenticazione dell'utente, che l'applicazione scambia con il
    token vero e proprio;
    \item \textbf{ID Token}: token in formato JWT~\cite{rfc7519}, composto da
    intestazione, \emph{payload} e firma. Il \emph{payload} contiene i
    \emph{claim} sull'utente, tra cui l'emittente (\texttt{iss}, ovvero
    \texttt{https://appleid.apple.com}), il destinatario (\texttt{aud}, cioè
    l'identificativo dell'applicazione), la scadenza (\texttt{exp}), l'identificativo
    univoco dell'utente (\texttt{sub}) e l'indirizzo email con il relativo stato
    di verifica (\texttt{email}, \texttt{email\_verified});
    \item \textbf{Firma digitale}: apposta da Apple con la propria chiave
    privata e verificabile da chiunque tramite le chiavi pubbliche che Apple
    pubblica.
\end{itemize}

L'applicazione non conosce la password dell'utente: la prova dell'identità è
costituita unicamente dal token. Nel caso in esame, tutte le credenziali
dell'attaccante erano di fatto sostituite da un token autentico ottenuto con un
parametro manipolato.

\subsection{Debolezze nella gestione e nella verifica dei token}
La debolezza principale risiedeva nell'\textbf{emissione} del token da parte
dell'IdP, non nella sua struttura crittografica: la firma era valida perché
generata realmente da Apple. Ne deriva che una verifica della sola firma da parte
delle applicazioni non era sufficiente a rilevare l'abuso~\cite{jain-apple-2020}.

Dal lato delle applicazioni, le indicazioni ufficiali prevedono di controllare
firma, emittente, destinatario, scadenza e identificativo dell'utente
(\texttt{sub})~\cite{apple-verifying-user}. I servizi che, oltre alla firma, si
affidavano solo all'indirizzo email come chiave di identificazione dell'account
erano i più esposti. Ne risulta un caso di \emph{fiducia eccessiva} nell'IdP:
quando il provider è compromesso o difettoso, le applicazioni prive di ulteriori
controlli ereditano direttamente la vulnerabilità~\cite{nist800-63c}.

\subsection{Il ruolo dell'email e dell'identità dell'utente nel processo di autenticazione}
L'indirizzo email ha un ruolo centrale perché molte applicazioni lo utilizzano
come identificativo per associare l'accesso federato a un account già
esistente. Questa scelta è comoda, ma implica che chi riesce a ottenere un token
con l'email di un altro utente venga riconosciuto come tale. L'email è infatti un
attributo \emph{descrittivo} dell'identità, che può cambiare nel tempo e che,
nel caso di Apple, può anche essere un indirizzo anonimo generato dal
servizio.

L'identificativo più affidabile è invece il claim \texttt{sub}, stabile e
univoco per ciascuna coppia utente--applicazione, che Apple indica come
riferimento per riconoscere l'utente~\cite{apple-verifying-user,openid-connect-core}.
Il caso del 2020 mostra quindi che la sicurezza dell'autenticazione federata
dipende dalla corretta associazione tra l'utente realmente autenticato e le
informazioni di identità riportate nel token, e che le applicazioni dovrebbero
basarsi su identificativi robusti e su ulteriori controlli, e non sulla sola
email.
